\Needspace{8\baselineskip}
\Article{292}
Die Zerlegungen von Zahlen in drei Quadrate unterscheiden wir (wie
oben bei den binären Formen) von den Darstellungen durch die Form
$xx+yy+zz$ so, dass wir bei jenen nur die Größe der Quadrate, bei diesen
aber auch ihre Reihenfolge und die Vorzeichen ihrer Wurzeln betrachten.
Wir halten also die Darstellungen $x=a$, $y=b$, $z=c$ und $x=a'$,
$y=b'$, $z=c'$ für verschieden, sofern nicht zugleich $a=a'$, $b=b'$,
$c=c'$ gilt. Die Zerlegungen in $aa+bb+cc$ und in $a'a'+b'b'+c'c'$ gelten dagegen
als eine, wenn diese Quadrate jenen ohne Rücksicht auf die Reihenfolge
gleich sind. Daraus folgt:

I.\ Eine Zerlegung der Zahl $M$ in die Quadrate $aa+bb+cc$ entspricht
$48$ Darstellungen, wenn keines $=0$ und alle ungleich sind;
aber $24$, wenn \textit{entweder} eines $=0$ und die übrigen
ungleich sind, \textit{oder} keines $=0$ ist und zwei gleich sind.
Sind dagegen in einer Zerlegung der gegebenen Zahl in drei Quadrate
zwei davon $=0$, \textit{oder} eines $=0$ und die übrigen
gleich, \textit{oder} alle gleich, so entspricht sie $6$,
\textit{oder} $12$, \textit{oder} $8$ Darstellungen. Dies kann
aber nur in den besonderen Fällen $M=1$, $2$ bzw.\ $3$
vorkommen, sofern die Darstellungen eigentlich sein sollen. Diese
Fälle seien ausgeschlossen. Die Anzahl aller Zerlegungen der Zahl
$M$ in drei Quadrate (ohne gemeinsamen Teiler) sei $E$, und
darunter seien $e$, in denen ein Quadrat $0$ ist, sowie
$e'$, in denen zwei Quadrate gleich sind. Jene kann man auch als
Zerlegungen in zwei Quadrate, diese als Zerlegungen in ein Quadrat
und das Doppelte eines Quadrats ansehen. Dann ist die Anzahl aller
eigentlichen Darstellungen der Zahl $M$ durch $xx+yy+zz$
\[
=24(e+e')+48(E-e-e')=48E-24(e+e')
\]
\SourceLeaf{351}
Aus der Theorie der binären Formen folgt aber leicht, dass $e$
entweder $=0$ oder $=2^{\mu-1}$ ist, je nachdem $-1$ quadratischer
Nichtrest oder Rest von $M$ ist; ebenso $e'=0$ oder $=2^{\mu-1}$,
je nachdem $-2$ Nichtrest oder Rest von $M$ ist. Dabei bezeichnet
$\mu$ die Anzahl der (ungeraden) Primfaktoren von $M$ (vgl.\
Art.\ 182; eine ausführlichere Darstellung lassen wir hier weg).
Daraus folgt leicht:
\[
\begin{aligned}
E&=2^{\mu-2}k,&&\text{wenn $-1$ und $-2$ beide $N.R.$ von $M$ sind;}\\
E&=2^{\mu-2}(k+2),&&\text{wenn beide Zahlen Reste sind; schließlich}\\
E&=2^{\mu-2}(k+1),&&\text{wenn eine Zahl Rest, die andere Nichtrest ist.}
\end{aligned}
\]
In den ausgeschlossenen Fällen $M=1$ und $M=2$ würde diese
Formel $E=\tfrac{3}{4}$ ergeben, während $E=1$ sein muss; für $M=3$ ergibt
sich dagegen richtig $E=1$, da sich die Ausnahmen gegenseitig ausgleichen.

Ist also $M$ eine Primzahl, so ist $\mu=1$, also $E=\tfrac{1}{2}(k+2)$, wenn
$M\equiv1\ (\mathrm{mod.}\ 8)$; $E=\tfrac{1}{2}(k+1)$, wenn $M\equiv3$ oder $\equiv5$. Diese besonderen Sätze
hat der berühmte Le Gendre durch Induktion gefunden und in seiner
schon öfter angeführten ausgezeichneten Abhandlung \textit{Hist.
de l'Ac. de Paris} 1785, S.\ 530 ff.\ mitgeteilt, wenn auch in etwas
anderer Form. Der Grund liegt vor allem darin, dass er die eigentliche
von der uneigentlichen Äquivalenz nicht unterschieden und daher
entgegengesetzte Klassen zusammengefasst hat.

II.\ Um alle Zerlegungen der Zahl $M$ in drei Quadrate (ohne
gemeinsamen Teiler) zu finden, braucht man nicht sämtliche eigentlichen
Darstellungen aller Formen $\varphi$, $\varphi'$, $\varphi''$ zu bestimmen.
Zunächst bestätigt man leicht, dass alle $(48)$ Darstellungen der
Form $\varphi$, die zu demselben Wert des Ausdrucks $\surd{}-(p,-q,r)$ gehören
(wobei $\varphi=(p,q,r)$ gesetzt wird), dieselbe Zerlegung der Zahl $M$
ergeben. Daher genügt es, eine davon zu haben, oder, was auf dasselbe
hinausläuft, nur alle verschiedenen Zerlegungen\HistoricalNote{*)}{Dabei
ist stets \textit{eigentliche} mitzudenken, wenn man diesen Ausdruck
von den Darstellungen auf die Zerlegungen übertragen will.} der
Form $\varphi$ in drei Quadrate aufzuschreiben; ebenso bei den übrigen
$\varphi'$, $\varphi''$ usw.\ Gehört ferner $\varphi$ zu einer nichtambigen
Klasse, so darf man die aus der entgegengesetzten Klasse gewählte
Form ganz übergehen; von zwei entgegengesetzten Klassen genügt es
also, eine einzige zu betrachten. Da die Wahl der Form aus jeder
Klasse ganz beliebig ist, nehmen wir an, aus der zur Klasse von
$\varphi$ entgegengesetzten Klasse werde die zu $\varphi$ entgegengesetzte
Form gewählt, und diese sei $=\varphi'$. Dann sieht man ohne Mühe:
Werden
\SourceLeaf{352}
die eigentlichen Zerlegungen der Form \(\varphi\) allgemein durch
\[
(gt+hu)^2+(g't+h'u)^2+(g''t+h''u)^2
\]
gegeben, so werden alle Zerlegungen der Form \(\varphi'\) durch
\[
(gt-hu)^2+(g't-h'u)^2+(g''t-h''u)^2
\]
ausgedrückt, und aus diesen folgen dieselben Zerlegungen der Zahl
\(M\) wie aus jenen. Ist schließlich \(\varphi\) eine Form aus einer
ambigen Klasse, die aber weder der Hauptklasse angehört noch zu
der Form \((2,0,\frac12 M)\) oder \((2,1,\frac12(M+1))\) äquivalent ist (je nachdem \(M\)
gerade oder ungerade ist), so darf man von den Werten des Ausdrucks
\(\surd-(p,-q,r)\) die Hälfte weglassen. Der Kürze wegen erklären wir diese
Abkürzung hier aber nicht näher. --- Dieselben Abkürzungen lassen
sich auch verwenden, wenn sämtliche eigentlichen Darstellungen
von \(M\) durch \(xx+yy+zz\) gesucht sind, da sich diese sehr leicht
aus den Zerlegungen ergeben.

Als Beispiel suchen wir alle Zerlegungen der Zahl \(770\) in drei
Quadrate. Hier ist \(\mu=3\), \(e=e'=0\), also \(E=2k\). Aus der Klassifikation
der positiven binären Formen mit der Determinante \(-770\), die
wir der Kürze wegen weglassen, da sie jeder nach Art.\ 231 leicht
aufstellen kann, ergibt sich die Anzahl der positiven Klassen
\(=32\). Alle sind eigentlich primitiv und auf \(8\) Geschlechter
verteilt, sodass \(k=4\) und daher \(E=8\) ist. Das Geschlecht mit
der charakteristischen Zahl \(-1\) muss bezüglich der Zahlen
\(5,7,11\) offenbar die besonderen Charaktere \(R5;\ N7;\ N11\) haben. Daraus
folgt nach Art.\ 263 leicht, dass sein Charakter bezüglich der
Zahl \(8\) gleich \(1\text{ und }3,\ 8\) sein muss. Nun finden sich in dem
Geschlecht mit dem Charakter \(1\text{ und }3,\ 8;\ R5;\ N7;\ N11\) vier Klassen, als deren Vertreter
wir die Formen \((6,2,129)\), \((6,-2,129)\), \((19,3,41)\), \(19,-3,41)\) wählen. Die
zweite und vierte Klasse lassen wir aber als entgegengesetzte der
ersten und dritten weg. Vier Zerlegungen der Form \((19,3,41)\) haben
wir schon in Art.\ 289 angegeben; daraus folgen die Zerlegungen
der Zahl \(770\) in \(9+361+400\), \(16+25+729\), \(81+400+289\), \(576+169+25\). Ebenso
findet man vier Zerlegungen der Form \(6tt+4tu+129uu\) in
\[
\begin{gathered}
(t-8u)^2+(2t+u)^2+(t+8u)^2,\\
(t-10u)^2+(2t+5u)^2+(t+2u)^2\\
(2t-5u)^2+(t+10u)^2+(t+2u)^2,\\
(2t+7u)^2+(t-8u)^2+(t-4u)^2
\end{gathered}
\]
\SourceLeaf{353}
die der Reihe nach aus den folgenden Werten des Ausdrucks \(\surd-(6,-2,129)\)
entstehen: \((48,369)\), \((62,-149)\), \((92,-159)\), \((202,61)\). Daraus folgen
die Zerlegungen der Zahl \(770\) in \(225+256+289\), \(1+144+625\), \(64+81+625\),
\(16+225+529\). Außer diesen acht Zerlegungen gibt es keine weiteren.

Was die Zerlegungen von Zahlen in drei Quadrate mit gemeinsamen
Teilern betrifft, so folgt dies so leicht aus der allgemeinen Theorie
in Art.\ 281, dass wir uns dabei nicht aufzuhalten brauchen.

\Needspace{8\baselineskip}
\begin{center}\textit{Beweis der Sätze von Fermat, dass sich jede ganze Zahl in drei Dreieckszahlen oder vier Quadrate zerlegen lässt.}\end{center}
\Article{293}
Die vorigen Untersuchungen liefern auch einen Beweis des berühmten
Satzes, dass sich \textit{jede positive ganze Zahl in drei Dreieckszahlen
zerlegen lässt}, den Fermat schon vor langer Zeit gefunden hat,
für den aber bisher ein strenger Beweis fehlte. Offenbar ergibt
jede Zerlegung der Zahl \(M\) in Dreieckszahlen
\[
\tfrac12 x(x+1)+\tfrac12 y(y+1)+\tfrac12 z(z+1)
\]
eine Zerlegung der Zahl \(8M+3\) in drei ungerade Quadrate
\[
(2x+1)^2+(2y+1)^2+(2z+1)^2
\]
und umgekehrt. Nach der vorigen Theorie lässt sich aber jede positive
ganze Zahl \(8M+3\) in drei Quadrate zerlegen, die notwendig ungerade
sind (vgl.\ die Anmerkung zu Art.\ 291). Die Anzahl der Zerlegungen
hängt sowohl von der Anzahl der Primfaktoren von \(8M+3\) als auch
von der Anzahl der Klassen ab, in die sich die binären Formen mit
der Determinante \(-(8M+3)\) verteilen. Ebenso viele Zerlegungen der
Zahl \(M\) in drei Dreieckszahlen gibt es. Wir setzen dabei voraus,
dass \(\frac12 x(x+1)\) für jeden ganzzahligen Wert von \(x\) als Dreieckszahl
angesehen wird. Wollte man die Null lieber ausschließen, müsste
man den Satz so ändern: Jede positive ganze Zahl ist entweder selbst
eine Dreieckszahl oder lässt sich in zwei oder drei Dreieckszahlen
zerlegen. Eine ähnliche Änderung wäre beim folgenden Satz nötig,
wenn man die Null von den Quadraten ausschließen wollte.

Aus denselben Grundsätzen folgt ein anderer Satz von Fermat: \textit{Jede
positive ganze Zahl lässt sich in vier Quadrate zerlegen.} Zieht
man von einer Zahl der Form \(4n+2\) ein beliebiges Quadrat (das
kleiner als diese Zahl ist), von einer Zahl der Form
\SourceLeaf{354}
\(4n+1\) ein gerades Quadrat und von einer Zahl der Form \(4n+3\)
ein ungerades Quadrat ab, so lässt sich der Rest in allen diesen
Fällen in drei Quadrate zerlegen, die gegebene Zahl also in vier.
Schließlich kann man eine Zahl der Form \(4n\) als \(4^{\mu}N\) darstellen,
wobei \(N\) zu einer der drei vorigen Formen gehört. Mit der
Zerlegung von \(N\) in vier Quadrate ist auch \(4^{\mu}N\) zerlegt.
Von einer Zahl der Form \(8n+3\) kann man auch ein Quadrat mit
durch vier teilbarer Wurzel abziehen, von einer Zahl der Form
\(8n+7\) ein Quadrat, dessen Wurzel das Doppelte einer ungeraden
Zahl ist, und von einer Zahl der Form \(8n+4\) ein ungerades Quadrat;
der Rest lässt sich dann in drei Quadrate zerlegen. Übrigens hatte
der berühmte La Grange diesen Satz schon bewiesen, \textit{Nouv.
Mém. de l'Ac. de Berlin} 1770, S.\ 123. Diesen von unserem ganz
verschiedenen Beweis hat der berühmte Euler in \textit{Actis Ac.
Petr. Vol.}\ II, S.\ 48 ausführlicher erklärt. --- Andere Sätze
Fermats, die gewissermaßen die Fortsetzung dieser Sätze bilden,
nämlich dass sich jede ganze Zahl in fünf Fünfeckszahlen, sechs
Sechseckszahlen, sieben Siebeneckszahlen \textit{usw.}\ zerlegen
lässt, sind bisher unbewiesen und scheinen tiefer liegende Grundsätze
zu erfordern.
